\section{Resultados Esperados}

A implementação da solução proposta tem como objetivo gerar ganhos concretos tanto no curto quanto no longo prazo, atuando diretamente na redução de riscos, aumento de eficiência operacional e evolução da arquitetura de integrações.

\subsection{Resultados no curto prazo}

No contexto imediato, espera-se que a solução:
\begin{itemize}
    \item Reduza significativamente o risco associado à migração das integrações do Rest para o Integrador;
    \item Permita validar o comportamento do novo sistema com base em dados reais de produção;
    \item Aumente a confiança da equipe técnica no processo de migração;
    \item Diminua a dependência de testes manuais e validações pontuais;
    \item Apoie decisões de rollout de forma mais segura e orientada por evidências.
\end{itemize}

\subsection{Resultados no médio prazo}

Com a evolução da solução, espera-se:
\begin{itemize}
    \item Aceleração do processo de migração das integrações ainda existentes no Rest;
    \item Redução do tempo necessário para integração de novas locadoras e drivers;
    \item Melhoria na manutenção de integrações já existentes;
    \item Redução de incidentes relacionados a falhas de integração;
    \item Maior previsibilidade na evolução do sistema.
\end{itemize}

\subsection{Resultados no longo prazo}

No longo prazo, a solução tende a se consolidar como um componente estratégico da arquitetura, possibilitando:
\begin{itemize}
    \item Evolução para uma plataforma contínua de observabilidade das integrações;
    \item Visibilidade consolidada sobre o comportamento das locadoras;
    \item Monitoramento de métricas técnicas e de negócio relevantes;
    \item Apoio à tomada de decisão com base em dados estruturados;
    \item Redução do custo operacional associado à manutenção e diagnóstico de problemas.
\end{itemize}

\subsection{Impacto organizacional}

Além dos ganhos técnicos, a solução contribui para a organização como um todo ao:
\begin{itemize}
    \item Promover uma abordagem orientada a dados na tomada de decisão;
    \item Reduzir o retrabalho e a necessidade de intervenções manuais;
    \item Melhorar a colaboração entre áreas técnicas e de negócio;
    \item Aumentar a maturidade do processo de integração com fornecedores.
\end{itemize}

\subsection{Conclusão}

De forma consolidada, a solução proposta permite transformar o processo de migração de integrações de um modelo baseado em validações manuais e baixa previsibilidade para um modelo estruturado, automatizado e orientado por evidências.

Como resultado, a organização passa a operar com maior segurança, eficiência e capacidade de evolução contínua, reduzindo riscos operacionais e aumentando a confiabilidade das integrações com locadoras.